\documentclass[11pt]{article}
\usepackage{amsmath,amssymb,amsthm}

\newtheorem{theorem}{Theorem}

\begin{document}

\begin{theorem}[Menelaus's Theorem]
\label{thm:menelaus}
Let $\triangle ABC$ be a triangle, and let a line $l$ meet the lines $BC$,
$CA$, and $AB$ (or their extensions) at points $D$, $E$, and $F$,
respectively, with $D,E,F$ distinct from $A,B,C$. Using signed lengths,
Menelaus's theorem states that
\[
\frac{\overline{AF}}{\overline{FB}}\cdot
\frac{\overline{BD}}{\overline{DC}}\cdot
\frac{\overline{CE}}{\overline{EA}}
=
-1.
\]
If one takes absolute values instead, then the corresponding unsigned product
is $1$.
\end{theorem}

\begin{proof}
An affine transformation preserves collinearity and preserves ratios of
directed segments on any fixed line. Therefore it is enough to prove the
theorem after sending
\[
A=(0,0),\qquad B=(1,0),\qquad C=(0,1).
\]
In these coordinates, the three sideline intersections can be written as
\[
F=(f,0),\qquad E=(0,e),\qquad D=(d,1-d)
\]
for some real numbers $f,e,d$.

Since $D,E,F$ are collinear, their coordinate matrix has determinant zero:
\[
\det
\begin{pmatrix}
f & 0 & 1 \\
0 & e & 1 \\
d & 1-d & 1
\end{pmatrix}
=0.
\]
Expanding the determinant gives
\[
-de+df+ef-f=0,
\]
or equivalently
\[
d(e-f)=f(1-e).
\]
Hence
\[
d=\frac{f(e-1)}{e-f},
\qquad
1-d=\frac{e(1-f)}{e-f}.
\]

Now compute the directed ratios along the three sidelines:
\[
\frac{\overline{AF}}{\overline{FB}}=\frac{f}{1-f},
\qquad
\frac{\overline{BD}}{\overline{DC}}=\frac{1-d}{d},
\qquad
\frac{\overline{CE}}{\overline{EA}}=\frac{1-e}{e}.
\]
Therefore
\[
\frac{\overline{AF}}{\overline{FB}}\cdot
\frac{\overline{BD}}{\overline{DC}}\cdot
\frac{\overline{CE}}{\overline{EA}}
=
\frac{f}{1-f}\cdot
\frac{e(1-f)}{f(e-1)}\cdot
\frac{1-e}{e}
=
-1.
\]
This proves Menelaus's theorem. Taking absolute values immediately gives the
unsigned version with product $1$.
\end{proof}

\end{document}
